\section{The Vibe: A Unit of Experience}

The tools are in hand. Now we put the one ingredient onto the crystal and the model begins. That ingredient is the vibe. This chapter says exactly what a vibe is and, just as important, what it is not.

\subsection{One kind of thing}

In this theory there is exactly one kind of entity. Not particles plus fields plus space plus minds. One kind of thing, the vibe, and everything else is a pattern of vibes. This is the strict version of an old idea called monism: reality is made of a single substance. Most monisms pick dead matter as the one substance and then cannot explain feeling. This one picks feeling.

A vibe is a unit of experience. The smallest possible dab of ``what it is like to be something.'' Not a particle that has an experience attached. The experience itself, however faint, at the smallest scale. When you reach the bottom of reality in this theory, you do not find lifeless dust. You find a tiny quiver of feeling.

\subsection{Why start with feeling}

The reason for this strange-sounding starting point is the hard problem we met in chapter 2. If you start with dead matter, you can never explain how feeling appears, because no arrangement of things that do not feel obviously adds up to something that does. The gap is total. So this theory refuses to open that gap in the first place. It puts feeling at the ground floor, where it cannot be mysteriously produced because it was never absent. The thing left to explain is then the reverse: not how feeling arises from matter, but how solid, lawful matter arises from feeling. That is a question the theory can actually answer, and the rest of the walkthrough answers it.

\subsection{A vibe is a cell}

Concretely, in the crystal we have built, each cell is one vibe. The vast hyperbolic graph of cells is a vast web of vibes. When two cells touch, two vibes are in direct contact. The whole machinery of Part II, the curved space, the tiling, the addressing, the automaton, is the machinery of how vibes are arranged and how they affect one another. The geometry is the structure of the field of feeling.

So you can now reread the crystal images with new eyes. That glowing heptagrid is not a picture of abstract tiles. It is a picture of feeling, cell by cell, each one a small experiencer, the whole forming one connected field.

\subsection{What a vibe is not}

Three quick guards against natural misreadings.

A vibe is not a little person or a little mind. It does not think or perceive a world. It is far simpler than that, a single bare quality, more like one pixel of feeling than a whole experience. Rich minds, as we will see, are huge organized crowds of vibes, not single ones.

A vibe is not made of anything more basic. It is the bottom. You do not get to ask what a vibe is built from, the way you cannot ask what the number one is built from. It is a starting element.

And a vibe is not located in a pre-existing space. There is no stage it sits on. Its only ``where'' is which other vibes it touches. Space, remember, is just the pattern of touching. The vibe comes first, the space is the shadow of how vibes connect.

\subsection{The single act}

A vibe does exactly one thing: it experiences other vibes, and is experienced by them. That mutual experiencing is the only activity in the entire theory. Everything, all of physics, all of mind, is built out of vibes experiencing vibes, over and over, on the growing crystal. The next chapters give that act its precise form: the felt charge a vibe carries (its tone), the relation of experiencing (the note), and the rule by which a vibe updates from what it experiences.

\textbf{Takeaway:} a vibe is the one and only kind of thing in the theory, a single unit of bare feeling, the cell of the crystal. Feeling is put at the ground floor so it need not be mysteriously produced. A vibe is not a mind, not made of parts, and not sitting in space. Its only act is to experience and be experienced.
